\documentclass[pdf,unicode,12pt,a4paper,oneside,fleqn]{article}

\usepackage{iftex}
\ifPDFTeX
  \usepackage[utf8]{inputenc}
  \usepackage[T2A]{fontenc}
  \usepackage{tempora}
\else
  \usepackage{fontspec}
  \setmainfont{texgyretermes-regular.otf}[
    BoldFont=texgyretermes-bold.otf,
    ItalicFont=texgyretermes-italic.otf,
    BoldItalicFont=texgyretermes-bolditalic.otf]
\fi
\usepackage[english,russian]{babel}
\usepackage{geometry}
\geometry{a4paper,left=2.54cm,right=2.54cm,top=2.54cm,bottom=2.54cm,
          headheight=0mm,headsep=0mm,footskip=0mm}
\usepackage{titlesec}
\usepackage{enumitem}
\usepackage{csquotes}
\usepackage{hyperref}
\hypersetup{hidelinks}

\titleformat*{\section}{\large\bfseries}
\titlespacing\section{0pt}{12pt plus 4pt minus 2pt}{0pt plus 2pt minus 2pt}
\setcounter{secnumdepth}{0}
\pagestyle{empty}
\frenchspacing
\parindent=0pt
\parskip=3pt

\newcommand{\StudentName}{Яковлев Артём Сергеевич}
\newcommand{\StudentGroup}{М8О-203БВ-25}
\newcommand{\TeacherName}{Миронов Евгений Сергеевич}

\begin{document}

\begin{titlepage}
\begin{center}
\large
МОСКОВСКИЙ АВИАЦИОННЫЙ ИНСТИТУТ\\
(НАЦИОНАЛЬНЫЙ ИССЛЕДОВАТЕЛЬСКИЙ УНИВЕРСИТЕТ)

\vspace{20pt}
Институт №8 <<Компьютерные науки и прикладная математика>>

Кафедра 806 <<Вычислительная математика и программирование>>
\end{center}

\vspace{60pt}
\begin{center}
\bfseries\large
Лабораторная работа №2\\
по курсу \enquote{Операционные системы}

\vspace{54pt}
Процессы и каналы\\
Вариант 22
\end{center}

\vfill
\begin{flushright}
\large
\begin{tabular}{rl}
Выполнил: & \StudentName \\
Группа: & \StudentGroup \\
Преподаватель: & \TeacherName \\
\end{tabular}
\end{flushright}

\vspace{92pt}
\begin{center}\large Москва, 2026\end{center}
\end{titlepage}
\pagebreak

\section{Цели и условие}
Цель работы -- научиться запускать дочерние процессы и передавать им данные через
каналы\cite{task}.

Программа сначала получает от пользователя имена двух файлов. Затем родительский
процесс запускает двух детей. Каждому ребёнку он может отправлять строки через
отдельный канал. Ребёнок переворачивает полученную строку и записывает результат
в свой файл.

По варианту 22 родитель выбирает первого ребёнка с вероятностью 80\%, а второго --
с вероятностью 20\%. Программа должна
работать на POSIX-системах и Windows.


\section{Программное и аппаратное обеспечение}
\begin{itemize}[noitemsep,topsep=0pt]
\item Язык программы -- C (стандарт C11); для сборки используется CMake.
\item Локальная проверка выполнена на macOS, компьютер ARM64, компилятор Apple Clang.
\item Сборка на Windows настроена в GitHub Actions на \texttt{windows-2022}.
\item Для проверки результата используется небольшой тест на Python.
\end{itemize}


\section{Метод решения}
Программа состоит из родительской программы \texttt{lab2\_parent} и дочерней
\texttt{lab2\_child}. Родитель запускает два экземпляра дочерней программы.
Для каждого он создаёт канал (pipe) и файл.

Работа программы проходит так:
\begin{enumerate}[noitemsep,topsep=0pt]
\item Родитель читает имена двух файлов и запускает обоих детей.
\item Для каждой новой строки он получает число от 0 до 99. Если число меньше 80,
строка отправляется первому ребёнку; иначе -- второму.
\item Ребёнок читает строку, переворачивает её и записывает в свой файл.
Например, \texttt{"Hello"} превращается в \texttt{"olleH"}.
\item Когда строки заканчиваются, родитель закрывает каналы и ждёт завершения детей.
\end{enumerate}

Длинные строки читаются с помощью увеличиваемого буфера. Для работы с процессами
используются разные функции POSIX и Windows, но порядок действий остаётся одним и тем же.
В Windows-реализации путь к программе и имена файлов должны состоять из
ASCII-символов; содержимое строк остаётся в UTF-8.

\section{Описание программы}
Основные файлы проекта:
\begin{itemize}[noitemsep,topsep=0pt]
\item \texttt{main.c} -- родитель: читает ввод, выбирает ребёнка и отправляет строки.
\item \texttt{child.c} -- ребёнок: переворачивает строки и выводит результат.
\item \texttt{line\_io.c} -- читает строки без заранее заданного ограничения длины.
\item \texttt{process\_posix.c} и \texttt{process\_win.c} -- создают каналы и запускают
процессы в соответствующей операционной системе.
\item \texttt{process.h} -- описывает структуру \texttt{ChildProcess}, в которой
родитель хранит данные о запущенном ребёнке и канале.
\item \texttt{CMakeLists.txt} -- выбирает нужный файл для текущей системы и собирает
обе программы.
\item \texttt{tests/test.py} проверяет результат работы для Windows.
\end{itemize}

Если при создании файла, канала или процесса происходит ошибка, программа сообщает
о ней и закрывает уже открытые ресурсы.


\section{Результаты}
На macOS программа прошла проверку. В одном запуске были введены строки
\texttt{Hello}, \enquote{Привет}, \texttt{abcdef} и \texttt{12345}.
В файле \texttt{first.txt} оказались \texttt{olleH}, \enquote{тевирП} и
\texttt{fedcba}; в \texttt{second.txt} -- \texttt{54321}.
Все строки записаны по одному разу и перевёрнуты правильно. В другом запуске
распределение по файлам может отличаться, потому что ребёнок выбирается случайно.

Опубликованная версия проекта находится в репозитории\cite{source}.
Её сборка и тест на Windows также прошли успешно в GitHub Actions\cite{ci}.



\section{Выводы}
В ходе работы написана программа, в которой родитель передаёт строки двум дочерним
процессам через отдельные каналы. Дети переворачивают строки и сохраняют их
в разные файлы. Такая схема подходит для задач, где поток независимых сообщений нужно распределить
между несколькими процессами. Например, каждый процесс может обрабатывать свою
часть входных данных и записывать результат в отдельный файл.

Основные трудности в реализации возникли при запуске
дочерних процессов и работе с каналами: важно передать ребёнку нужный конец
канала и вовремя закрыть ненужные концы. Иначе ребёнок может продолжать ждать
ввод после того, как родитель закончил отправлять строки. Для Windows также
потребовался отдельный код запуска процессов.


Проверка на macOS и Windows показала правильную работу на латинских и русских строках.


\begin{thebibliography}{99}
\bibitem{task} Лабораторная работа №2 по курсу \enquote{Операционные системы}.
Условие задания, вариант 22.

\bibitem{source} F14my. \enquote{OS-lab2}.
\href{https://github.com/F14my/OS-lab2}{Исходный код}.
\bibitem{ci} GitHub Actions. \enquote{Windows test}.
\href{https://github.com/F14my/OS-lab2/actions/runs/38629281921}{Успешный запуск для коммита 7193d5b}
\bibitem{microsoft}
Microsoft Learn. Creating a Child Process with Redirected Input and Output.
\href{https://learn.microsoft.com/en-us/windows/win32/procthread/creating-a-child-process-with-redirected-input-and-output}
{Документация Win32 API}.
\end{thebibliography}

\end{document}
